\section{Signed forcing and gap transport}
\label{sec:bounds}

The hierarchy and the complex-power formula describe the evolution
exactly. Their quantitative use is governed by the real forcing
\begin{equation}
 F=12Z(2)+16\delta Z(3),\qquad
 m_\sigma=2m^2-F,\qquad \delta_\sigma=2-4m\delta,
 \label{eq:consumer-system}
\end{equation}
where $m=\ell^2M$, $\delta=D/\ell^2$ and
$\sigma=(t-t_0)/\ell^2$. A lower bound for $m$ requires an
\emph{upper} bound for $F$. Reality of $F$ alone does not provide it.
We first compare the direct and contour estimates at the same information,
then derive the conditional transport bound. The comparison and static
example extend \cite{BakerGR05} as proved in \cite{BakerCM06}.

\subsection{Absolute estimates and preservation of joint signs}

The elementary shell count in \eqref{eq:count-safe} gives, for $n\ge1$,
\begin{equation}
 \mathcal A_n:=\sum_j|\mu_j|^n\le B_nL,
 \qquad B_n=14\Cu\gamma^n,\qquad \gamma=q_0/a^2.
 \label{eq:integer-bound}
\end{equation}
For completeness, the band $|\Re w|<1$ contains at most $4\Cu L$
occurrences, and the shell $k\le|\Re w|<k+1$ contains at most
$2\Cu(2L+\log(k+1))$. The kernel there is bounded by
$q_0\ell^2/k^2$. Using
$\sum k^{-2n}\le2$, $\sum k^{-2n}\log(k+1)\le3$ and $a\le L$
bounds the far shells by $10\Cu\gamma^nL$ and proves
\eqref{eq:integer-bound}.
Since $|\mu_j|\le\gamma$, the direct estimate is
\begin{equation}
 |F|\le (12\gamma+16d\gamma^2)\mathcal A_1\le C_FL,
 \qquad C_F=12B_2+16dB_3.
 \label{eq:direct-forcing}
\end{equation}

Fix the damping $\epsilon>0$ of Section~\ref{sec:cauchy}. Write
$C_\epsilon=(12+16d)\e^{4\epsilon}$ and let $\Theta$ be its branch
argument bound. On the two verticals $q=1,4$,
\[
 |K_\epsilon(q+\ii v,\delta)Z(q+\ii v)|
 \le C_\epsilon\mathcal A_1
        \frac{\e^{-\epsilon v^2+\Theta|v|}}{\sqrt{1+v^2}}.
\]
The absolute two-line envelope is therefore bounded by
\begin{equation}
 \frac1{2\pi}\int_\R
       (|K_\epsilon Z(4+\ii v)|+|K_\epsilon Z(1+\ii v)|)\dd v
 \le C_{\rm env}\mathcal A_1,
 \quad C_{\rm env}=\frac{2C_\epsilon}{\pi}
       \sqrt{\frac\pi\epsilon}\e^{\Theta^2/(4\epsilon)}.
 \label{eq:contour-envelope}
\end{equation}
Completing the square proves this estimate. At fixed collar parameters
$a,d$ and fixed damping $\epsilon$, its height order is the same as
\eqref{eq:direct-forcing}, and its displayed coefficient is
larger: minimizing $\e^{4\epsilon}/\sqrt\epsilon$ gives
$C_{\rm env}/(12+16d)\ge2\sqrt{8\e/\pi}>1$, whereas
$12\gamma+16d\gamma^2<12+16d$.
This comparison concerns these explicit envelopes, not every possible
signed contour estimate.

\begin{proposition}[Comparison for a fixed grouping]\label{prop:group-bound}
Partition the exterior occurrences into prescribed finite groups $G$,
and retain the assigned logarithms. Put
\[
 Z_G(p)=\sum_{j\in G}\mu_j^p,\quad
 I_{q,G}(v)=K_\epsilon(q+\ii v,\delta)Z_G(q+\ii v),\quad
 F_G=\sum_{j\in G}(12\mu_j^2+16\delta\mu_j^3),
\]
where the power notation uses those logarithms. Then
\begin{equation}
 F_G=\frac1{2\pi}\int_\R(I_{4,G}-I_{1,G})\dd v,
 \qquad
 |F_G|\le\frac1{2\pi}\int_\R|I_{4,G}-I_{1,G}|\dd v
 \le\frac1{2\pi}\int_\R(|I_{4,G}|+|I_{1,G}|)\dd v.
 \label{eq:group-bound}
\end{equation}
All group sums and integrals in these estimates interchange absolutely.
Thus estimating each group by an absolute contour norm costs at least
$\sum_G|F_G|$, the direct signed-residue norm for the same grouping.
For groups closed under conjugation as occurrence multisets, with compatible
logs, the same statement holds for the folded
norm $\pi^{-1}\int_0^\infty|\Re(I_{4,G}-I_{1,G})|\dd v$.
\end{proposition}

\begin{proof}
Apply the positive-rectangle residue formula to the finite entire
sum $Z_G$. The horizontal edges tend to zero by Gaussian damping.
The residues at $2$ and $3$ are independent of the assigned logarithms.
The triangle inequalities give \eqref{eq:group-bound}. On either
vertical, the sum over groups of the absolute integrands is bounded
by the integrable majorant preceding \eqref{eq:contour-envelope}.
The corresponding horizontal group sum is bounded by a constant times
$\e^{-\epsilon Y^2+\Theta Y}/Y$. This proves all complete-divisor
passages. For a compatible group, conjugation changes $v$ to $-v$,
giving the folded expression and the same triangle inequality.
\end{proof}

Proposition~\ref{prop:group-bound} leaves joint signs between groups
available on either side of the identity. Discarding them in the direct
formula but retaining them in the contour formula would compare
different information. The sharp zero-cutoff residue error in
Section~\ref{sec:cauchy} likewise belongs to both evaluations.

\subsection{A complete static example with linear forcing}

The following construction tests what the slice hypotheses can force.
It changes the canonical divisor and is not an example of its heat
evolution. The explicit constants ensure that the class is fixed as
$T$ grows.

\begin{proposition}[Scope of a static forcing bound]\label{prop:static}
There is a family of real-even entire functions of order one with a
counted double target, collar $a=4,d=1$, width-one strip, growing zero
count, unit-count and static safe-line bounds with fixed enlarged
constants, whose complete complex-power sums are holomorphic on
$\Re p>1/2$, have compatible Schwarz symmetry and admit the fixed-damping
contours of Section~\ref{sec:cauchy}, but satisfy
\begin{equation}
 m=-\frac1{32}+O(L^{-1}),\qquad \frac FL\longrightarrow\frac3{256}>0.
 \label{eq:static-asymptotic}
\end{equation}
In particular, this static class admits no uniform upper bound $F=o(L)$,
even on the subclass with $m\le-1/64$.
\end{proposition}

\begin{proof}
Take $c=T=\e^L$, $L\ge16$, $D=0$, $b=8$ and $J=\lfloor L\rfloor$.
For this separate slice construction use the inherited strip, unit-count
and safe-line bounds of the canonical field at time zero. The background
bounds hold on the closed slab $[0,1/2]$ independently of whether the
particular pair tube $I$ contains zero \cite{BakerBase,BakerGR05}.
Remove from $H(0,z)$ every occurrence in the closed bands
$|\Re\rho\mp c|\le3$ by division by their complete monic polynomial.
Insert at $+c$ the double target and $J$ occurrences at each
displacement $\pm b\ell$, and $J+1$ at each displacement
$\pm\ii b\ell$. Reflect all insertions at $-c$. The division is
entire, and conjugation and negation invariance make the resulting
function real and even. Finite polynomial multiplication and division
preserve the order bound; the retained growing count forces order one.
The strip and the required collar hold, with the remaining background
at horizontal distance greater than $3$.

The inherited canonical growing count and its logarithmic discrepancy
are the static inputs described in \cite[Section 3]{BakerGR05}.
There are at most $12\Cu L$ removed occurrences per band and at most
$4J+4\le5L$ insertions per band. Hence the unit-count constant can
be enlarged to $\Cu+5$, and the cumulative-count discrepancy constant
by $24\Cu+10$. The logarithmic derivative changes by a signed sum
of at most $(24\Cu+10)L$ reciprocal distances. On the line $v+4\ii$,
if $|v|\ge c/2$ use vertical distance at least $3$ and
$\log(2+|v|)\ge L/2$; otherwise use horizontal distance at least
$c/3$. The safe-line constant can be enlarged to
$\Cu+5(24\Cu+10)$, uniformly in $c$.
The corresponding lower safe line follows by conjugation.

For the near core assign conjugate logarithms $\pm\pi$ to the negative
values of $Q_j$ and real logarithms to the positive ones. Its complete
contribution is
\begin{equation}
 Z_{\rm core}(p)=b^{-2p}\{2J+2(J+1)\cos(\pi p)\}.
 \label{eq:static-core}
\end{equation}
All remaining occurrences are retained in a remainder $R$. The retained
background and every opposite-band insertion, including the reflected
double, satisfy $|\Re w|>3$, $|\Im w|\le1$. Hence
$\Re(w^2)>8$, and we assign principal logs on the entire noncore.
These are conjugation compatible and have argument magnitude less than
$\pi/2$. Unit-shell counting, including the opposite insertions and
opposite target, gives
for $1\le u\le r<\infty$
\begin{equation}
 |R(u+\ii v)|\le (10\Cu+5)L^{1-2u}\e^{\pi|v|/2}.
 \label{eq:static-tail}
\end{equation}
Indeed the physical absolute power sum of exponent $2u$ over the
retained background is at most $10\Cu L$; the opposite insertions
cost at most $5Lc^{-2u}$. Multiplication by $\ell^{2u}$ proves the
bound. The same shell argument for any $u>1/2$ proves normal
convergence on compact subsets of that half-plane. Compatible logs
give Schwarz symmetry, and Gaussian damping pays the complete vertical
integrals and vanishing horizontal edges. Evaluating
\eqref{eq:static-core}--\eqref{eq:static-tail} at $1$ and $2$ gives
\[
 m=-2/b^2+R(1),\qquad
 F=12(4J+2)/b^4+12R(2),
\]
which proves \eqref{eq:static-asymptotic} and the stated exclusion.
\end{proof}

The same construction retains the inherited broad transverse upper
bounds with fixed enlarged constants \cite[Section 3]{BakerGR05}.
If their original coefficient is $B_\beta$ for $0<\beta<1/2$, the
added nonreal occurrences lie
at height $8/L$ and number at most $3L$ per band. For every dyadic
interval $[U,2U]$ meeting the positive inserted nonreal projection,
$c\in[U,2U]$, including endpoint cuts, one has
$L/2\le\log U\le2L$. The comparison
$U\log U\,\e^{-\beta v\log U}\ge4L\e^{-16\beta}$ whenever
$v<8/L$ allows coefficient $B_\beta+\e^{16\beta}$; deletion only
reduces the counts. This is still a slice construction. It supplies
neither the unchanged theta source nor a common independently established
canonical time tube. In particular Proposition~\ref{prop:static} does
not exclude a signed estimate using the full canonical time law.

\subsection{A conditional Riccati comparison}

\begin{theorem}[Forcing upper bound and transport cost]\label{thm:transport}
Suppose $F(\sigma)\le Q$ on the entire established scaled tube
$0\le\sigma\le\Sigma=(t_1-t_0)/\ell^2$, for a real constant $Q$
that may depend on $T$.
Set $A_0=B_1L$, $K_+=\max(Q,0)$ and $h=\sqrt{K_+/2}$. Define
\begin{equation}
 \begin{array}{ll}
 f(\sigma)=\cosh(2h\sigma)+(A_0/h)\sinh(2h\sigma),&
 g_{\rm tr}(\sigma)=\sinh(2h\sigma)/(2h),\quad h>0,\\[2pt]
 f(\sigma)=1+2A_0\sigma,&g_{\rm tr}(\sigma)=\sigma,\quad h=0.
 \end{array}
 \label{eq:comparison-functions}
\end{equation}
Then, with primes here denoting scaled-time derivatives,
\begin{equation}
 m(\sigma)\ge-\frac{f'(\sigma)}{2f(\sigma)},\qquad
 m(\sigma)\ge-h-\frac1{2\sigma}\quad(\sigma>0).
 \label{eq:riccati-barrier}
\end{equation}
For $0\le x\le y\le\Sigma$, put
$W(x,y)=\exp(-4\int_x^y m(u)\dd u)$. Then
\begin{equation}
 \begin{split}
 W(x,y)&\le[f(y)/f(x)]^2,\\
 \int_x^y W(u,y)\dd u&\le\frac{f(y)}{f(x)}g_{\rm tr}(y-x),\\
 W(x,y)\delta(x)&=\delta(y)-2\int_x^y W(u,y)\dd u.
 \end{split}
 \label{eq:transport-cost}
\end{equation}
These bounds retain the finite initial layer at $\sigma=0$.
\end{theorem}

\begin{proof}
The absolute count gives $m(0)\ge-A_0$, and
$m'\ge2m^2-K_+$. The function $v=-f'/2f$ solves
$v'=2v^2-K_+$ with $v(0)=-A_0$. Since
$(m-v)'-2(m+v)(m-v)\ge0$, an integrating factor gives $m\ge v$.
For $h>0$, $f'/2f\le h\coth(2h\sigma)\le h+1/(2\sigma)$;
for $h=0$ use the displayed linear $f$. Integrating the first bound
gives the multiplier bound in \eqref{eq:transport-cost}.

Both $f$ and $g_{\rm tr}$ solve $u''=2K_+u$. Their Wronskian,
with the second solution translated by $x$, is
\[
 f(u)g_{\rm tr}'(u-x)-f'(u)g_{\rm tr}(u-x)=f(x).
\]
Consequently integrating $(g_{\rm tr}(u-x)/f(u))'$ gives
\[
 f(y)^2\int_x^y\frac{\dd u}{f(u)^2}
        =\frac{f(y)}{f(x)}g_{\rm tr}(y-x),
\]
which proves the cost bound. Finally
$\partial_u[W(u,y)\delta(u)]=2W(u,y)$ by
\eqref{eq:consumer-system}. Integrating proves the exact last identity.
\end{proof}

An independent terminal lower anchor
$\delta(y)\ge2f(y)g_{\rm tr}(y-x)/f(x)$ suffices for
$\delta(x)\ge0$. Equality with this upper allowance guarantees
nonnegativity, including $x=0$. Equality with the exact cost
$2\int_x^y W(u,y)\dd u$ gives $\delta(x)=0$, whereas strict
separation of the terminal value from that actual cost gives strict
positivity. The forward version uses an independent incoming
$\delta(x)$ and yields
\[
 |\delta(y)|\le[f(y)/f(x)]^2|\delta(x)|
                 +2f(y)g_{\rm tr}(y-x)/f(x).
\]
In physical variables $W$ integrates $M$ over physical time, and the
inhomogeneous cost for $D$ carries the additional factor $\ell^2$.
Neither the forcing estimate nor the initial moment bound supplies
either gap anchor or the tube itself.

Using only \eqref{eq:direct-forcing}, fixed collar parameters give
$K_+=C_FL$, an $O(\sqrt L)$ interior negative allowance after a
fixed positive scaled buffer, and an
$O(L)\exp(O(\sqrt L))$ upper multiplier on a fixed scaled duration.
The contour representations proved here do not improve this allowance
by the absolute estimates in Proposition~\ref{prop:group-bound}.
